\documentclass[a4paper,12pt]{article}
\usepackage[utf8]{inputenc}
\usepackage[T1]{fontenc}
\usepackage{lmodern}
\usepackage{geometry}
\geometry{a4paper, margin=1in}
\usepackage{graphicx}
\usepackage{listings}
\usepackage{xcolor}
\usepackage{hyperref}
\usepackage{fancyhdr}

\pagestyle{fancy}
\fancyhf{}
\fancyhead[L]{Massimo Mantineo}
\fancyhead[C]{Hands-On 18}
\fancyhead[R]{A.A. 2024/2025}
\fancyfoot[C]{\thepage}

\title{Relazione Hands-On 18: MapReduce Distribuito con Ray.io}
\author{Massimo Mantineo (Matricola: 541924)}
\date{Aprile 2026}

\definecolor{codegreen}{rgb}{0,0.6,0}
\definecolor{codegray}{rgb}{0.5,0.5,0.5}
\definecolor{codepurple}{rgb}{0.58,0,0.82}
\definecolor{backcolour}{rgb}{0.95,0.95,0.92}

\lstdefinestyle{mystyle}{
    backgroundcolor=\color{backcolour},   
    commentstyle=\color{codegreen},
    keywordstyle=\color{magenta},
    numberstyle=\tiny\color{codegray},
    stringstyle=\color{codepurple},
    basicstyle=\ttfamily\footnotesize,
    breakatwhitespace=false,         
    breaklines=true,                 
    captionpos=b,                    
    keepspaces=true,                 
    numbers=left,                    
    numbersep=5pt,                  
    showspaces=false,                
    showstringspaces=false,
    showtabs=false,                  
    tabsize=2
}
\lstset{style=mystyle}

\begin{document}

\maketitle
\tableofcontents
\newpage

\section{Introduzione}
L'Hands-On 18 affronta lo studio e l'implementazione di un sistema di calcolo distribuito basato sul paradigma \textbf{MapReduce}. A differenza degli esercizi precedenti svolti in C e focalizzati sulla comunicazione inter-processo a basso livello o socket, questo esperimento utilizza \textbf{Python} e il framework \textbf{Ray.io} per gestire la parallelizzazione e la distribuzione dei task su cluster (reali o virtuali).

L'obiettivo specifico è il conteggio delle frequenze dei termini presenti nel testo del celebre romanzo \textit{"Il Signore degli Anelli"}, analizzando come i task vengono suddivisi, eseguiti in parallelo e infine aggregati.

\section{Stato dell'Arte}
Il paradigma MapReduce, reso celebre da Google, si basa su due fasi principali:
\begin{itemize}
    \item \textbf{Map}: I dati di input vengono suddivisi in piccoli pezzi e processati in parallelo da diversi worker.
    \item \textbf{Reduce}: I risultati parziali dei worker vengono aggregati per produrre il risultato finale.
\end{itemize}

Framework moderni come \textbf{Ray} permettono di implementare tali pattern in modo efficiente, offrendo astrazioni semplici per l'esecuzione remota di funzioni e la gestione delle dipendenze, superando i limiti della libreria standard multiprocess di Python.

\section{Metodologia e Implementazione}
L'architettura del sistema prevede:
\begin{enumerate}
    \item \textbf{Splitting}: Il testo di input (\texttt{lotr.txt}) viene suddiviso in $N$ blocchi.
    \item \textbf{Remote Map Tasks}: Ogni blocco è processato da una funzione Ray remota (\texttt{@ray.remote}) che produce un \texttt{Counter} locale.
    \item \textbf{Aggregation}: I risultati vengono scaricati dal cluster tramite \texttt{ray.get()} e fusi nel Counter globale.
\end{enumerate}

Il codice principale è contenuto in \texttt{mapreduce\_lotr.py}.

\section{Risultati Sperimentali}
L'esecuzione sul testo completo del romanzo ("The Lord of the Rings") ha prodotto i seguenti risultati (Top 10 parole per frequenza):

\begin{verbatim}
Starting distributed MapReduce on lotr.txt...
Report generated: lotr_report.txt
Top 10 words:
  the: 11705
  and: 7555
  of: 5083
  to: 3948
  a: 3728
  he: 3038
  in: 2961
  i: 2825
  it: 2780
  that: 2539
\end{verbatim}

I risultati mostrano la tipica distribuzione delle lingue naturali, dove articoli e congiunzioni dominano il conteggio. La distribuzione tramite Ray è avvenuta correttamente, inizializzando un cluster locale sulla macchina di test e completando l'analisi di migliaia di parole in pochi secondi.

\section{Conclusioni e Sviluppi Futuri}
L'esperimento dimostra l'efficacia di Ray per l'elaborazione distribuita. Sviluppi futuri potrebbero includere:
\begin{itemize}
    \item Implementazione di fasi di Reduce distribuite (Shuffling).
    \item Utilizzo di dataset massivi caricati da cloud storage (AWS S3/GCS).
    \item Filtraggio di "stop-words" per un'analisi semantica più accurata.
\end{itemize}

\end{document}
